\section*{Impact Statement}

The (partial) automation of intellectual labor will have significant societal and socioeconomic impact, both positive and negative.
One current hurdle is the automation of general interactive decision-making tasks, as our benchmark demonstrates.
If solved, this would lead to more autonomous AI systems and agents that could facilitate or even partly take over intellectual labor across a broad range of applications and domains.
We believe this vision warrants both enthusiasm and caution, and the scientific community and developers of AI technology bear responsibility in raising awareness of potential negative socioeconomic outcomes, and in helping develop potential mitigations. 
Having said that, our work is a benchmark that reports the current state of affairs (showing the limitations of current frontier LMs) and serves as one yard stick to measure future innovations and progress against, but does not propose innovations to improve capabilities of AI systems.
While our benchmark may contribute indirectly by helping develop more capable AI systems faster, we believe this is far outweighed by the benefit of having a clearer picture of current capabilities and progress provided by our benchmark.

